\section{Additional spectral estimates}\label{app:spectral}\label{app:metric}
The estimates in Section~\ref{sec:metric} reduce the size of a selected
minor to three bounds for a node Gram matrix. We first prove the
complementary singular-product bound and the Schur bound. We then give
finite estimates for the row weights and the three node bounds. All
parameters remain free, so the same estimates apply when the polynomial
space or the number of evaluation points changes.

\subsection{The complementary singular product}\label{supp:complementary}
The forward estimate uses an affine map taking the disk into itself.
For a finite polynomial space, an expanding affine map is also available.
To estimate its largest singular values, we pass to the inverse map and
estimate the complementary product.

\begin{proof}[Proof of Lemma~\ref{lem:complementary}]
The matrix $T_J$ is triangular, with diagonal entries $1,d,\ldots,d^{J-1}$.
Consequently
\[
 \det(T_J^*T_J)=d^{J(J-1)}.
\]
The product of its largest $m$ squared singular values is this determinant
times the product of the largest $k=J-m$ squared singular values of
$T_J^{-1}$. The inverse is composition by $(z-c_0)/d$.

Put $D_{s_0}=\operatorname{diag}(s_0^j)_{0\leq j<J}$, with $G_0,s_0$
as in the lemma. Then $T_J^{-1}D_{s_0}$ is composition by
\[
 \psi(z)=-\frac{c_0r_0}{G_0}+\frac{r_0}{G_0}z.
\]
This map is a strict disk contraction. Indeed, $G_0$ is the larger root
of
\[
 X^2-(1+r_0^2)X+(1-c_0^2)r_0^2=0.
\]
The left side at $X=(1+c_0)r_0$ equals
$-(1+c_0)r_0(1-r_0)^2<0$, so $G_0>(1+c_0)r_0$.
Changing the sign of the translation does not change singular values,
by the unitary reflection $f(z)\mapsto f(-z)$.
Lemma~\ref{lem:composition}, applied to $\psi$, gives squared singular
values $G_0^b r_0^{2l}$, $l\geq0$, for the infinite composition operator.
Compression to polynomials of degree less than $J$ cannot increase them.

For matrices $X,Y$, the norm inequality
$\|\bigwedge^k(XY)\|\leq
\|\bigwedge^kX\|\,\|\bigwedge^kY\|$ bounds the product of the largest
$k$ singular values of a product. Apply it to
\[
 T_J^{-1}=(T_J^{-1}D_{s_0})D_{s_0}^{-1}.
\]
The first factor contributes at most $G_0^{bk}r_0^{k(k-1)}$ to the
squared singular product. If $s_0\leq1$, the largest $k$ diagonal
entries of $D_{s_0}^{-1}$ have indices $J-k,\ldots,J-1$; if $s_0\geq1$,
they have indices $0,\ldots,k-1$. Their squared products are respectively
$s_0^{-k(2J-k-1)}$ and $s_0^{-k(k-1)}$. Multiplying these bounds and
taking logarithms gives \eqref{eq:complementary-cost}.
\end{proof}

\subsection{The Schur estimate}\label{app:schur-proof}
The Schur bound sums all principal minors of a given size. Its proof has
two parts. The Meixner factorization first compares each minor with a
gamma moment determinant. Schur polynomials describe how that determinant
depends on the selected orders, and the Cauchy identity sums this
dependence over all selections.

We use the following standard notation. A partition of length at most
$m$ is an integer sequence $\lambda_1\geq\cdots\geq\lambda_m\geq0$;
write $|\lambda|=\sum_i\lambda_i$. Its Schur polynomial is
\[
 s_\lambda(x_1,\ldots,x_m)
 =\frac{\det[x_j^{\lambda_i+m-i}]_{i,j=1}^m}
        {\det[x_j^{m-i}]_{i,j=1}^m}.
\]
The quotient extends to a symmetric polynomial when some variables
coincide. At $1^m=(1,\ldots,1)$ it has value
\begin{equation}\label{eq:schur-dimension}
 s_\lambda(1^m)
 =\prod_{1\leq i<j\leq m}
       \frac{\lambda_i-\lambda_j+j-i}{j-i}.
\end{equation}
We will use the Littlewood--Richardson product rule and the Cauchy
identity in their usual Schur specializations; see~\cite{fulton} and
\cite[Section~2]{forresterwarnaar}.

\begin{lemma}[Gamma moments]\label{lem:gamma-schur}
Let $b>0$ and $m\geq1$, let $0\leq j_0<\cdots<j_{m-1}$, and put
$D_T=\det[(b)_{j_i+j_k}]_{i,k=0}^{m-1}$. Then
\begin{equation}\label{eq:gamma-minor-bound}
 D_T\leq
 \left(\prod_{0\leq i<k<m}\frac{j_k-j_i}{k-i}\right)^2
 \prod_{i=0}^{m-1}i!(b)_{2j_i-i}.
\end{equation}
\end{lemma}
\begin{proof}
Let $\mathrm d\mu_b(x)=x^{b-1}e^{-x}\mathrm dx/\Gamma(b)$ on
$(0,\infty)$. Its $r$th moment is $(b)_r$. Expanding both determinants
and integrating term by term gives Andreief's identity in the form
\[
 D_T=\frac1{m!}\int_{(0,\infty)^m}
          \bigl(\det[x_k^{j_i}]_{i,k=0}^{m-1}\bigr)^2
          \prod_{k=0}^{m-1}\mathrm d\mu_b(x_k).
\]
Set $\lambda_{m-i}=j_i-i$. The alternant formula shows that
$\det[x_k^{j_i}]$ is the Vandermonde determinant times $s_\lambda(x)$,
up to a sign. For the consecutive orders $0,\ldots,m-1$, the moment
determinant is
\[
 D_0=\det[(b)_{i+k}]_{i,k=0}^{m-1}
     =\prod_{i=0}^{m-1}i!(b)_i.
\]
To see the last equality, factor $(b)_i$ from row $i$. The remaining
entry $(b+i)_k$ is monic of degree $k$ in $b+i$, so its determinant is
the Vandermonde product $\prod_{i=0}^{m-1}i!$.

Write $\Delta(x)=\prod_{i<k}(x_k-x_i)$. Normalize the density
$\Delta(x)^2\prod_k\mathrm d\mu_b(x_k)$ to have mass one, and denote
expectation for it by $\mathbb E$. The preceding calculation gives
$D_T=D_0\mathbb E[s_\lambda(X)^2]$. For any partition $\nu$ of length
at most $m$, a second use of Andreief gives
\begin{equation}\label{eq:gamma-schur-average}
 \mathbb E[s_\nu(X)]
 =s_\nu(1^m)\prod_{i=1}^m(b+m-i)_{\nu_i}.
\end{equation}
To evaluate it, Andreief writes the expectation as
\[
 \mathbb E[s_\nu(X)]
 =\frac{\det[(b)_{\nu_i+m-i+m-j}]_{i,j=1}^m}{D_0}.
\]
Factor $(b)_{\nu_i+m-i}$ from row $i$. The remaining entries are the
monic polynomials $(b+\nu_i+m-i)_{m-j}$. Dividing their Vandermonde
determinant by the one for $\nu=0$ gives
$s_\nu(1^m)$ by \eqref{eq:schur-dimension}; the ratios of the row
factors give the rising factorials in \eqref{eq:gamma-schur-average}.

The Littlewood--Richardson rule writes
$s_\lambda^2=\sum_\nu c_{\lambda\lambda}^{\nu}s_\nu$, with
nonnegative coefficients. Every occurring $\nu$ has size $2|\lambda|$
and is dominated by $2\lambda$: for every $r$,
$\sum_{i\leq r}\nu_i\leq2\sum_{i\leq r}\lambda_i$.
Terms with more than $m$ parts vanish in $m$ variables.
The decreasing vector $(b+m-i+2\lambda_i)_{i=1}^m$ therefore
majorizes $(b+m-i+\nu_i)_{i=1}^m$. Convexity of $\log\Gamma$ yields
\[
 \prod_{i=1}^m(b+m-i)_{\nu_i}
 \leq\prod_{i=1}^m(b+m-i)_{2\lambda_i}.
\]
Apply \eqref{eq:gamma-schur-average} to the product expansion and
evaluate that expansion at $1^m$. We obtain
\[
 D_T\leq D_0s_\lambda(1^m)^2
                    \prod_{i=1}^m(b+m-i)_{2\lambda_i}.
\]
Substituting $\lambda_{m-i}=j_i-i$ and using
\eqref{eq:schur-dimension} gives \eqref{eq:gamma-minor-bound}.
\end{proof}

\begin{proof}[Proof of \eqref{eq:schur}]
The case $m=0$ is immediate, so assume $m\geq1$. Use the notation
$M_{jk}(b,w)$, $\mathcal B$, $\sigma_0$, $c_*$ and $z_*$ from
Section~\ref{sec:metric}. The Meixner identities
\eqref{eq:meixner-norm}--\eqref{eq:meixner-expansion} give
\begin{equation}\label{eq:positive-meixner}
 M_{jk}(b,w)
 =(1-w)^{-b-j-k}(b)_j(b)_k
 \sum_{l=0}^{\min(j,k)}
 \frac{\fall jl\fall kl}{l!(b)_l}w^{j+k-l}.
\end{equation}
These are finite consequences of the generating identity
\eqref{eq:nb-generating}; the standard Meixner formulas are also recorded
in~\cite[Sections~18.19 and~18.23]{dlmf}.

Fix $T=\{j_0<\cdots<j_{m-1}\}$, write $s=\sum_{j\in T}j$, and put
$V_{lj}=\fall jl/\sqrt{l!(b)_l}$ for $0\leq l,j<J$.
Cauchy--Binet applied to \eqref{eq:positive-meixner} gives
\begin{equation}\label{eq:meixner-minor-positive}
 (1-w)^{bm+2s}\det M[T,T]
 =\prod_{j\in T}(b)_j^2
   \sum_{\substack{L\subseteq\{0,\ldots,J-1\}\\|L|=m}}
      \det(V_{L,T})^2w^{2s-\sum_{l\in L}l}.
\end{equation}
A nonzero term matches each $l\in L$ to an order $j\geq l$.
Thus $\sum_{l\in L}l\leq s$. After taking out $w^s$, the remaining
polynomial has nonnegative coefficients and is at most its value at
$w=1$. The finite Vandermonde identity is
\[
 (b)_j(b)_k\sum_l\frac{\fall jl\fall kl}{l!(b)_l}=(b)_{j+k}.
\]
For example, it follows by comparing the coefficient of $x^j$ in
$(1-x)^{-b-k}=(1-x)^{-b}(1+x/(1-x))^k$.
Consequently \eqref{eq:meixner-minor-positive}, with the column and
mass factors restored, gives
\begin{equation}\label{eq:schur-gamma-reduction}
 \det G[T,T]\leq\mathcal B^m(1-w)^{-bm}c_*^{2s}D_T.
\end{equation}

Write $a_*=J-m$ and $s_*=ma_*+m(m-1)/2$, and set
$d_i=a_*+i-j_i$ for $0\leq i<m$. The inequalities on the $j_i$ imply
\[
 a_*\geq d_0\geq\cdots\geq d_{m-1}\geq0.
\]
Thus $d$ is a partition, $s=s_*-|d|$, and the dimension factor in
\eqref{eq:gamma-minor-bound} is $s_d(1^m)$.
The ratio $(b)_{2j_i-i}/(b)_{2a_*+i}$ removes $2d_i$ factors, all
at least $b$. Combining Lemma~\ref{lem:gamma-schur} with
\eqref{eq:schur-gamma-reduction} therefore gives
\[
 \frac{\det G[T,T]}
 {\mathcal B^m(1-w)^{-bm}c_*^{2s_*}
                         \prod_{i=0}^{m-1}i!(b)_{2a_*+i}}
 \leq s_d(1^m)^2 z_*^{|d|}.
\]
The assumption $bc_*>1$ says $0<z_*<1$. Sum over the order sets $T$
and enlarge the sum to all partitions of length at most $m$.
The Schur Cauchy identity, specialized to two lists of $m$ variables
all equal to $\sqrt{z_*}$, gives
\[
 \sum_d s_d(1^m)^2z_*^{|d|}=(1-z_*)^{-m^2}.
\]
This is precisely the remaining factor in \eqref{eq:schur}.
\end{proof}

\subsection{The signed row sum}
We now estimate the exact formulas at finite $R$ and $S$. Throughout,
$R=KS$, $1\leq J\leq R$, and the weights are
\eqref{eq:row-weights}, with $a>0$, $0<v<1$, and $\beta<0$.
The estimates below keep the row sum separate, since it can be negative.

Write $L(x)=x\log x$, with $L(0)=0$, and define
\[
 f(x)=L(x)-L(a+x)+L(a)-x\log v,
 \qquad C_v=-a\log(1-v).
\]
\begin{lemma}\label{lem:spectral-row-finite}
Suppose $aR>1$ and $f(1)<0$. Then
\begin{equation}\label{eq:row-finite}
 \frac1W\sum_\Lambda\lambda_{r,s}\leq
 3\int_0^1(1-x)f(x)\,\mathrm dx+\frac{9C_v}{R}
 +\frac\beta{2K}(1-S^{-2})
 +\frac{3(1+S^{-1})}{2R}\delta_R,
\end{equation}
where
\[
 \delta_R=\log(R+1)+1+\log\frac{a+1}{a-1/R}.
\]
\end{lemma}
\begin{proof}
Integral comparison gives, for $r\geq0$ and $b>1$,
\[
 \log(r!)\leq r\log r-r+\log(r+1)+1,
 \qquad
 \log((b)_r)\geq\int_0^r\log(b+t)\,\mathrm dt
                        -\log\frac{b+r}{b-1}.
\]
For the second inequality, compare the sum with the integral over
$[-1,r-1]$ and bound the two endpoint intervals. Substituting $b=aR$
in the row weights yields
$\lambda_{r,s}\leq Rf(r/R)+\beta s+\delta_R$.

The function $f$ is convex, $f(0)=0$, and its minimum on $[0,\infty)$
is $-C_v$. Thus $f\leq0$ on $[0,1]$ and its total variation there is
at most $2C_v$. The function $(1-x)f(x)$ has total variation at most
$3C_v$: apply the product bound for variation, using $|f|\leq C_v$.
There are $S-\lfloor r/K\rfloor\geq S(1-r/R)$ rows with first
exponent $r$. Since $f$ is nonpositive, their contribution is at most
$RS(1-r/R)f(r/R)$. The left Riemann sum differs from the integral
by at most $3C_v/R$, giving the first two terms of
\eqref{eq:row-finite} after division by $W=R^2S/3$.

Finally,
$\sum_\Lambda s=KS(S^2-1)/6$ and
$|\Lambda|=R(S+1)/2$. Summing $\beta s+\delta_R$ gives the last two
terms of \eqref{eq:row-finite}.
\end{proof}
The integral in \eqref{eq:row-finite} is elementary; it involves only
integrals of $x\log x$ and $x^2\log x$, after translation by $a$.

\subsection{Column capacities and the forward estimate}
Put $\varepsilon=J/R$. For a fixed upper squared radius $0<q\leq1$
and a column scale $\tau>0$, define
\[
 g_\tau(x)=L(x)+L(a+x)-L(a)
                  +2x(\ell-\log\tau-1),
 \qquad E_{\rm col}=\frac{\log(R+1)+1}{R}.
\]
The price \eqref{eq:column-prices} satisfies
\begin{equation}\label{eq:column-finite-error}
 \frac{\nu_j}{R}\leq g_\tau(j/R)+E_{\rm col}.
\end{equation}
Indeed, use the factorial upper bound above and
$\log((aR)_j)\leq\int_0^j\log(aR+x)\,\mathrm dx$.
The terms containing $\log R$ cancel those in the column scale, and
$j+1\leq J\leq R$ bounds the factorial remainder.

The sum of the largest $m$ column prices can be bounded without deciding
which orders occur. Partition $[0,\varepsilon]$ into $N_z$ equal bins.
For a bin $I=[u,v]$, give it the upper cost
$\gamma_I=\max\{0,g_\tau(u),g_\tau(v)\}$ and capacity
\[
 c_z=\frac{\varepsilon}{N_z}+\frac1R.
\]
The function $g_\tau$ is convex, so $\gamma_I$ bounds it throughout
$I$. Each bin contains at most $Rc_z$ of the possible points $j/R$.
Assign boundary points to just one neighboring bin, and define
\[
 C_{q,\tau}(t)
 =\max\left\{\sum_I\gamma_I u_I:
                0\leq u_I\leq c_z,\quad\sum_Iu_I=t\right\}.
\]
This function is nondecreasing and piecewise linear. Filling bins in
decreasing order of $\gamma_I$ evaluates it. Every selection of $m$
orders gives a feasible choice at $t=m/R$; hence
\[
 \frac{\nu_m^*(\tau)}{R^2}\leq C_{q,\tau}(t)+tE_{\rm col}.
\]
The replacement of negative bin costs by zero ensures monotonicity and
only enlarges the upper bound.

Retain $\delta_q,v_q,w,\mathcal B$ from
\eqref{eq:moment-matrix}. For the forward choice $0<\xi<1$, take
$\Gamma,\tau$ as in Section~\ref{sec:metric}, and set
\[
 d_{q,\xi}=a\log\frac\Gamma{1-w}+\frac{\delta_q}{K}
             +\frac{\log R}{R}+E_{\rm col}-\frac{\log\xi}{R}.
\]
Since $\log\mathcal B/R\leq\log R/R+\delta_q/K$, the forward
bound \eqref{eq:forward-cost} gives
\begin{equation}\label{eq:forward-finite}
 \frac{F_q^{\rm f}(m)}{R^2}
 \leq C_{q,\tau}(t)+d_{q,\xi}t+t^2\log\xi.
\end{equation}
The last term in $d_{q,\xi}$ retains the difference between
$m(m-1)$ and $m^2$.

To use \eqref{eq:H-occupancy}, we need a bound throughout an interval
of possible values of $t$, rather than only at its endpoints. On a cell
$[l,h]\subseteq[0,\varepsilon]$, a bound for
$F_q^{\rm f}(m)/R^2-a_0t$ is
\begin{equation}\label{eq:forward-cell}
 C_{q,\tau}(h)
 +\max\{(d_{q,\xi}-a_0)l,(d_{q,\xi}-a_0)h\}
 +l^2\log\xi.
\end{equation}
The three terms follow respectively from the monotonicity of
$C_{q,\tau}$, the sign of the linear coefficient, and $\log\xi<0$.

\subsection{Finite form of the complementary estimate}
For $\kappa>0$, put
\[
 d=\frac{\kappa}{1+\sqrt w},\qquad c_0=d\sqrt w,
 \qquad \tau=\frac{d(1-w)}{\sqrt{v_q}}.
\]
These are compatible with the affine map \eqref{eq:affine-map}.
Choose $0<r_0<1$ and use $G_0,s_0$ from
Lemma~\ref{lem:complementary}. Define
\begin{align*}
 c_1&=-a\log G_0-\varepsilon
                  (2\log r_0-\log s_0-|\log s_0|),\\
 c_2&=\log r_0-|\log s_0|<0,\\
 e_q&=-a\log(1-w)+\frac{\delta_q}{K}
          +\frac{\log R}{R}+E_{\rm col}
          +\frac{\max(0,-\log d)}{R}.
\end{align*}
The complementary bound for the node minor, divided by $R^2$, is at most
\begin{equation}\label{eq:complement-finite}
 C_{q,\tau}(t)+\varepsilon(a+\varepsilon)\log G_0
                          +(c_1+e_q)t+c_2t^2.
\end{equation}
To verify the finite term, substitute $J=\varepsilon R$, $m=tR$ and
$k=(\varepsilon-t)R$ into \eqref{eq:complementary-cost}. Using
$\log s_0=\log d+\log r_0-\log G_0$, both branches become
\[
 \varepsilon(a+\varepsilon)\log G_0+c_1t+c_2t^2
       -\frac{t\log d+(\varepsilon-t)\log G_0}{R}.
\]
Here $G_0\geq1$ and $0\leq t\leq\varepsilon$, so the last term is at
most $t\max(0,-\log d)/R$. Restoring the mass and column prices gives
\eqref{eq:complement-finite}.

On a cell $[l,h]$, subtract $a_0t$ and use
\begin{equation}\label{eq:complement-cell}
 \begin{split}
 C_{q,\tau}(h)+\varepsilon(a+\varepsilon)\log G_0
 &+\max\{(c_1+e_q-a_0)l,(c_1+e_q-a_0)h\}\\
 &+c_2l^2.
 \end{split}
\end{equation}
As in the forward estimate, every term bounds the whole cell.

\subsection{The finite error in the Schur estimate}
At a positive upper squared radius $q$, put
\[
 \zeta=\ell+\tfrac12\log v_q-\log(1-w),\qquad
 z_*=a^{-2}e^{-2\zeta},\qquad
 I(x)=\tfrac12x^2\log x-\tfrac14x^2,
\]
with $I(0)=0$. In \eqref{eq:schur}, the parameters of the completed
Gram matrix satisfy $c_*=e^\zeta/R$ and $bc_*=ae^\zeta$.
Define
\begin{align}
 g_\varepsilon(t)={}&t\bigl(\delta_q/K-a\log(1-w)\bigr)
          +(2\varepsilon t-t^2)(\zeta-1)+I(t)\notag\\
 &+I(a+2\varepsilon-t)-I(a+2\varepsilon-2t)
               -ta\log a-t^2\log(1-z_*).
 \label{eq:schur-finite}
\end{align}

\begin{proposition}\label{prop:schur-finite-error}
Suppose $R>3$, $0<a\leq1$, $0<\varepsilon=J/R\leq1$,
$\varepsilon(a+2\varepsilon)\leq1$, and $ae^\zeta>1$.
For $0\leq m\leq J$ and $t=m/R$, the logarithm of the right side of
\eqref{eq:schur}, divided by $R^2$, is at most
\begin{equation}\label{eq:schur-error}
 g_\varepsilon(t)+E_0,\qquad
 E_0=\frac{\varepsilon}{R}
           (2\log R+2-\log\varepsilon-\log a).
\end{equation}
\end{proposition}
\begin{proof}
The assertion is immediate for $m=0$. For $m\geq1$, put $a_*=J-m$,
as in \eqref{eq:schur}. In the product there, use
\[
 \log(i!)\leq i\log i-i+\log(i+1)+1,\qquad
 \log((aR)_n)\leq\int_0^n\log(aR+x)\,\mathrm dx.
\]
For $i=0,\ldots,m-1$, set $x_i=i/R$ and take $n=2a_*+i$.
The terms containing $\log R$ have total coefficient
$\sum_i(2a_*+2i)=2s_*$, and cancel those in
$c_*^{2s_*}=\exp(2s_*(\zeta-\log R))$.
The remaining sum, divided by $R^2$, is the left Riemann sum
$R^{-1}\sum_{i=0}^{m-1}f_t(x_i)$, where
\[
 \begin{split}
 f_t(x)={}&x\log x+(a+2\varepsilon-2t+x)
                         \log(a+2\varepsilon-2t+x)\\
          &-a\log a-2\varepsilon+2t-2x.
 \end{split}
\]

The other finite terms have simple upper bounds. Since $a\leq1$ and
$ae^\zeta>1$, we have $\zeta>0$. Thus replacing
$2s_*/R^2=2\varepsilon t-t^2-t/R$ by $2\varepsilon t-t^2$ drops the
nonpositive term $-t\zeta/R$. The factor $\mathcal B^m$ contributes
$t\delta_q/K+t\log S/R$, whose second term is at most
$\varepsilon\log R/R$. Since $i+1\leq J\leq R$, the factorial
remainders contribute at most $\varepsilon(\log R+1)/R$.

It remains to estimate the Riemann sum. For $0<x\leq t$,
\[
 f_t'(x)=\log\bigl(x(a+2\varepsilon-2t+x)\bigr)\leq0,
\]
because $x\leq t\leq\varepsilon$ and
$\varepsilon(a+2\varepsilon)\leq1$. Its total variation is at most
\[
 -\int_0^t(\log x+\log a)\,\mathrm dx
 =t(1-\log t-\log a)
 \leq\varepsilon(1-\log\varepsilon-\log a).
\]
The last inequality follows by differentiation, since $at\leq1$.
For a decreasing function, the left Riemann sum exceeds the integral
by at most its variation divided by $R$. Finally, direct integration
gives
\[
 \begin{split}
 \int_0^t f_t(x)\,\mathrm dx
 ={}&I(t)+I(a+2\varepsilon-t)-I(a+2\varepsilon-2t)\\
    &-ta\log a-2\varepsilon t+t^2.
 \end{split}
\]
Combining this integral with the mass, $c_*^{2s_*}$, and
$(1-z_*)^{-m^2}$ gives \eqref{eq:schur-finite}. The three error bounds
sum to $E_0$, proving \eqref{eq:schur-error}.
\end{proof}

To obtain a bound on a whole cell, we use concavity. Differentiating
\eqref{eq:schur-finite} twice gives, for $t>0$,
\[
 g_\varepsilon''(t)=
 \log\frac{t(a+2\varepsilon-t)}
 {e^{2\zeta}(a+2\varepsilon-2t)^4(1-z_*)^2}.
\]
In particular, the sufficient condition
\begin{equation}\label{eq:schur-concavity}
 \varepsilon(a+\varepsilon)z_*<a^2(1-z_*)^2
\end{equation}
makes $g_\varepsilon$ concave on $[0,\varepsilon]$. Indeed, the
numerator is at most $\varepsilon(a+\varepsilon)$, the fourth power
in the denominator is at least $a^4$, and
$e^{2\zeta}=1/(a^2z_*)$.
For a cell with midpoint $t_0>0$ and half-width $h_0$, its tangent
therefore gives the upper bound
\begin{equation}\label{eq:schur-cell}
 g_\varepsilon(t_0)-a_0t_0
        +h_0|g_\varepsilon'(t_0)-a_0|+E_0
\end{equation}
for $g_\varepsilon(t)-a_0t+E_0$ throughout the cell. This also applies
to the first cell, whose left endpoint is zero, by continuity.

\subsection{Combining the finite estimates}
Partition $[0,\varepsilon]$ into $N_m$ cells. At each positive upper
radius $q$, the forward and complementary formulas give bounds on every
cell. Include \eqref{eq:schur-cell} whenever the hypotheses of
Proposition~\ref{prop:schur-finite-error} and
\eqref{eq:schur-concavity} hold. For each cell, take the minimum over
the available bounds and their parameter choices. Let $\Phi(q)$ be
the maximum of these cell bounds and zero. Each chosen bound holds
throughout its cell. Taking $F_q(m)$ to be the minimum of the
corresponding exact node bounds, we obtain
\[
 F_q(m)\leq a_0Rm+R^2\Phi(q)\qquad(0\leq m\leq J).
\]
Thus $\Phi$ satisfies \eqref{eq:H-occupancy}. Radial monotonicity
(Lemma~\ref{lem:radial}) and the node count in
\eqref{eq:metric-mesh} now give the upper bound for the full minor.
The signed row sum is supplied by \eqref{eq:row-finite}.

For the finite comparisons, the additional parameter choices are
\[
 \kappa\in\{5/2,3,4,6,10\},
\]
\begin{multline*}
 \xi,r_0\in\{1/100,3/100,3/50,1/10,3/20,1/5,3/10,2/5,\\
                 1/2,3/5,7/10,4/5,9/10,19/20,99/100,9999/10000\}.
\end{multline*}
The estimates themselves apply to any finite collections satisfying
$\kappa>0$ and $0<\xi,r_0<1$. They may also use an upper bound for
$\ell$: each fixed principal minor is multiplied by
$e^{2\ell\sum_{j\in T}j}$, so increasing $\ell$ enlarges it.
All radial upper endpoints are positive; the original node minors at
$q=0$ are covered by Lemma~\ref{lem:radial}.
